\section{Triển khai}

\subsection{Danh sách linh kiện điện tử sử dụng trong sản phẩm}

Để bảo đảm tính khả thi và độ tin cậy trong điều kiện vận hành liên tục, nhóm đã lựa chọn các module phần cứng tiêu chuẩn và có sẵn trên thị trường. Bảng \ref{tab:danh_sach_linh_kien} tổng hợp danh mục linh kiện, số lượng, đơn giá và ID chức năng liên quan (ảnh bố trí thực tế xem tại Mục 2.1).

\begin{center}
\scriptsize
\renewcommand{\arraystretch}{0.95}
\begin{tabularx}{\textwidth}{|c|X|c|c|p{3.7cm}|c|}
\hline
\textbf{STT} & \textbf{Tên linh kiện} & \textbf{SL} & \textbf{Giá (VNĐ)} & \textbf{Link Shopee tham khảo} & \textbf{ID CN} \\ \hline
1 & Vi điều khiển ESP32 DevKit V1 (30 chân) & 1 & 84.114 & \href{https://shopee.vn/product/58532883/43882498412}{shopee.vn/esp32-devkit} & ID 11, 12 \\ \hline
2 & Cảm biến nhiệt độ \& độ ẩm DHT22 & 1 & 45.144 & \href{https://shopee.vn/product/16504852/10244874574}{shopee.vn/dht22-am2302} & ID 1 \\ \hline
3 & Module quang trở LM393 & 1 & 10.000 & \href{https://shopee.vn/product/484901669/46818176298}{shopee.vn/lm393-ldr} & ID 4 \\ \hline
4 & Còi báo chủ động Active Buzzer 5V & 1 & 6.500 & \href{https://shopee.vn/product/8004713/18650576346}{shopee.vn/active-buzzer-5v} & ID 3 \\ \hline
5 & Nút nhấn cơ 4 chân $6 \times 6\text{ mm}$ (Gói 50) & 1 gói & 19.000 & \href{https://shopee.vn/search?keyword=nut\%20nhan\%204\%20chan\%206x6mm}{shopee.vn/nut-nhan-6x6} & ID 6 \\ \hline
6 & Đèn LED đỏ 5mm (Combo 20 bóng) & 1 gói & 5.000 & \href{https://shopee.vn/product/1610645505/26194886231}{shopee.vn/led-do-5mm} & ID 6 \\ \hline
7 & Gói 50 điện trở vạch 1/4W & 1 gói & 3.000 & \href{https://shopee.vn/product/542242728/46614464788}{shopee.vn/dien-tro-1-4w} & Mạch phụ \\ \hline
8 & Màn hình LCD 1602 kèm I2C PCF8574 & 1 bộ & 25.000 & \href{https://shopee.vn/product/404478549/18549939018}{shopee.vn/lcd1602-i2c} & ID 5 \\ \hline
9 & Module đọc thẻ RFID RC522 (13.56 MHz) & 1 bộ & 25.000 & \href{https://shopee.vn/product/337676632/11610099363}{shopee.vn/rfid-rc522} & ID 2 \\ \hline
10 & Thẻ từ RFID Mifare Classic 1K S50 & Kèm bộ & Kèm bộ & \href{https://shopee.vn/product/337676632/11610099363}{Đi kèm module RC522} & ID 2 \\ \hline
11 & Breadboard MB-102 (Bo test 830 lỗ) & 1 & 25.000 & \href{https://shopee.vn/product/58532883/55663851300}{shopee.vn/breadboard-mb102} & Kết nối \\ \hline
12 & Dây cắm Dupont 20cm (40 sợi) & 1 tép & 15.000 & \href{https://shopee.vn/product/1071100439/47763941551}{shopee.vn/day-dupont-20cm} & Kết nối \\ \hline
13 & Hộp nhựa kỹ thuật ABS bảo vệ mạch & 1 & 21.000 & \href{https://shopee.vn/search?keyword=hop\%20nhua\%20ky\%20thuat\%20abs}{shopee.vn/hop-nhua-abs} & Vỏ hộp \\ \hline
14 & Nguồn Adapter 9V 1A / Cáp Micro USB & 1 & 25.000 & \href{https://shopee.vn/product/213519943/7319530329}{shopee.vn/adapter-9v-1a} & Nguồn điện \\ \hline
\multicolumn{3}{|r|}{\textbf{Tổng chi phí dự kiến:}} & \multicolumn{3}{l|}{\textbf{308.758 VNĐ} (\textit{Ba trăm lẻ tám nghìn bảy trăm năm mươi tám đồng})} \\ \hline
\end{tabularx}
\captionof{table}{Danh mục linh kiện điện tử và dự trù kinh phí thực hiện sản phẩm}
\label{tab:danh_sach_linh_kien}
\end{center}

\subsubsection{Cơ sở vật lý của các linh kiện và cảm biến sử dụng}

Nhằm gắn kết chặt chẽ phần cứng với kiến thức học phần \textit{Vật lý cho Công nghệ Thông tin}:
\begin{enumerate}[leftmargin=1.2cm, noitemsep, topsep=1pt]
    \item \textbf{DHT22 (Nhiệt ẩm):} Đo nhiệt độ qua nhiệt điện trở NTC ($R$ giảm theo hàm mũ khi $T$ tăng do electron vượt vùng cấm bán dẫn); đo độ ẩm dựa trên sự thay đổi điện dung $C = \frac{\varepsilon \varepsilon_0 S}{d}$ của lớp polymer khi hấp phụ hơi nước.
    \item \textbf{Module LM393 (Quang trở):} Dựa trên \textbf{hiệu ứng quang điện trong} của chất bán dẫn CdS; photon chiếu vào giải phóng hạt dẫn tự do làm điện trở giảm từ hàng trăm k$\Omega$ xuống vài trăm $\Omega$.
    \item \textbf{Module RFID RC522:} Dựa trên \textbf{hiện tượng cảm ứng điện từ cộng hưởng} (định luật Faraday) ở tần số 13.56 MHz, từ trường biến thiên nạp năng lượng cho chip trong thẻ phát ngược mã UID.
    \item \textbf{Đèn LED đỏ 5mm:} Dựa trên \textbf{hiện tượng phát quang do tái hợp} electron - lỗ trống khi phân cực thuận lớp tiếp giáp p-n, giải phóng photon ánh sáng đỏ đặc trưng ($\lambda \approx 620 - 750\text{ nm}$).
    \item \textbf{Active Buzzer 5V:} Hoạt động dựa trên \textbf{hiệu ứng áp điện nghịch} tạo dao động sóng âm nghe được ($f \approx 2.5\text{ kHz}$) trên màng rung cơ học khi có xung điện áp kích.
\end{enumerate}

\clearpage
\subsection{Kế hoạch thực hiện}

Kế hoạch triển khai được phân bổ trong 3 tuần (15/11 -- 05/12), bảo đảm mỗi chức năng chỉ do 1 thành viên chịu trách nhiệm chính theo đúng quy chế môn học:

\begin{center}
\scriptsize
\renewcommand{\arraystretch}{1.0}
\begin{tabularx}{\textwidth}{|c|X|p{2.6cm}|p{3.2cm}|}
\hline
\textbf{ID CN} & \textbf{Nội dung công việc thực hiện} & \textbf{Thời gian} & \textbf{Phân công**} \\ \hline
\textbf{ID 1} & Lập trình đọc cảm biến DHT22, kiểm tra checksum và xuất bản dữ liệu lên topic \texttt{study/dht} & 15/11 -- 18/11 (15h) & Trần Nguyễn Tuấn Minh \\ \hline
\textbf{ID 2} & Tích hợp đầu đọc RFID RC522 qua SPI, đọc UID và đối soát dữ liệu phiên học trên Supabase & 22/11 -- 25/11 (20h) & Nguyễn Quốc Dương \\ \hline
\textbf{ID 3} & Lập trình điều khiển còi Active Buzzer 5V theo nhịp Pomodoro và tạo nút bấm hẹn giờ trên web & 15/11 -- 19/11 (15h) & La Anh Hào \\ \hline
\textbf{ID 4} & Lập trình đọc ADC1 quang trở LM393, tính toán mức sáng và hiển thị cảnh báo trên Dashboard & 15/11 -- 19/11 (12h) & Nguyễn Quốc Dương \\ \hline
\textbf{ID 5} & Lập trình giao tiếp I2C điều khiển LCD 1602, giải mã chuỗi văn bản từ topic \texttt{study/lcd} & 22/11 -- 25/11 (15h) & La Anh Hào \\ \hline
\textbf{ID 6} & Lập trình ngắt ngoài nút nhấn GPIO27 (chống rung phím) và điều khiển đèn LED đỏ báo bận & 18/11 -- 21/11 (10h) & Nguyễn Quốc Dương \\ \hline
\textbf{ID 7} & Thiết kế CSDL Supabase PostgreSQL, cấu hình bảng \texttt{sensor\_data} và \texttt{study\_sessions} & 18/11 -- 21/11 (18h) & Trần Nguyễn Tuấn Minh \\ \hline
\textbf{ID 8} & Xây dựng giao diện Node-RED Dashboard hiển thị kim đo Gauge, đồ thị chuỗi thời gian và bảng lịch sử & 22/11 -- 27/11 (22h) & Trần Nguyễn Tuấn Minh \\ \hline
\textbf{ID 9} & Xây dựng luồng đếm thời gian liên tục và tích hợp Telegram Bot API phát tin nhắc nghỉ $>60$p & 25/11 -- 28/11 (16h) & Nguyễn Quốc Dương \\ \hline
\textbf{ID 10} & Xây dựng luồng tổng hợp dữ liệu cuối tuần và tích hợp Gmail SMTP gửi email báo cáo tiến độ & 29/11 -- 02/12 (16h) & Nguyễn Quốc Dương \\ \hline
\textbf{ID 11} & Xây dựng luồng xác thực tài khoản, mã hóa băm mật khẩu bcrypt và bảo mật dữ liệu cá nhân & 25/11 -- 28/11 (18h) & La Anh Hào \\ \hline
\textbf{ID 12} & Lập trình Web Server SoftAP Captive Portal trên ESP32 cho phép cấu hình Wi-Fi linh hoạt & 29/11 -- 02/12 (20h) & La Anh Hào \\ \hline
\textbf{ID 13} & Thu thập dữ liệu thử nghiệm và lập trình mô hình phân tích điều kiện học tập tối ưu trên Node-RED & 29/11 -- 03/12 (20h) & Trần Nguyễn Tuấn Minh \\ \hline
\textbf{Tích hợp} & Lắp ráp vỏ hộp 3D, thử nghiệm liên thông phần cứng với Node-RED và Supabase, đóng gói sản phẩm & 03/12 -- 05/12 (16h) & Cả nhóm (Minh, Dương, Hào) \\ \hline
\end{tabularx}
\captionof{table}{Kế hoạch triển khai chi tiết và phân công cá nhân cho 13 chức năng}
\label{tab:ke_hoach_thuc_hien}
\end{center}

\noindent \textit{\footnotesize ** Ghi chú: Mỗi chức năng (ID 1 đến ID 13) chỉ phân công cho duy nhất một thành viên chịu trách nhiệm chính theo đúng quy định đồ án cuối kỳ. Giai đoạn nước rút (03/12 -- 05/12) cả 3 thành viên phối hợp kiểm thử liên thông toàn diện và nghiệm thu.}
